%% =====================================================================
%%  B1-T-07-03 -- what B1 contributes to "The Three-Second Read"
%%  -------------------------------------------------------------------
%%  LAYER A: APPEND-ONLY. Written once, then left alone. Material found
%%  only here sits in the `unique` slot and is protected by rule R10.
%% =====================================================================
%% @kind: source
%% @id: B1-T-07-03
%% @book: B1
%% @topic: T-07-03
%% @locator: Tactic 6, pp. 72--74
%% @status: distilled
%% @unique: yes
%% @integrated: yes
%% @updated: 2026-09-19

\dossier{B1}{T-07-03}{\bookshortB1}{Tactic 6, pp. 72--74}

\begin{distilled}
  \item The key to a person's true thoughts is said to lie in the area around the eyes, where emotions are mirrored clearly.
  \item The procedure: look the subject straight in the eyes, ask the direct question, then focus on the immediate expression across the eyes and upper cheekbones.
  \item Within three seconds the gut feeling about the question is said to be readable almost always.
  \item Used this way the procedure is claimed to defeat attempts to mislead by lying, which is offered as the explanation for why it is nearly impossible to deceive an experienced deceiver.
  \item A worked interpersonal example is given: introducing the subject casually, then asking the direct question and reading the immediate reaction rather than the denial that follows.
\end{distilled}

\begin{unique}
  \item The three-second window, the specific facial region to observe, and the procedure of pairing a direct question with observation of the immediate reaction rather than of the answer.
  \item The reflexive claim that experienced deceivers are hard to deceive because they run this read continuously.
\end{unique}
